\question[5]
Liste ist unser selbstgebastelter Datentyp.
Das Bild zeigt die Situation der Liste a.

\includegraphics[height=3cm]{\pfad/eintraege_01.png}

a. Was liefern die folgenden Anweisungen: 
\begin{lstlisting}
print(a.elem()) 
print(a.endpos())
\end{lstlisting}
b. Zeichne die Situation nach der Ausführung von a.insert(12) \\
c. Was liefern jetzt die folgenden Anweisungen: 
\begin{lstlisting}
print(a.elem()) 
print(a.endpos())
\end{lstlisting}
d. Wir gehen nochmal von der Anfangssituation aus. Zeichne die Situation nach Ausführung von a.delete() \\
e. Was liefern jetzt die folgenden Anweisungen: 
\begin{lstlisting}
print(a.endpos())
print(a.elem())
\end{lstlisting}


\begin{solutionbox}{9.25cm}
a. 9, False \\
b. 

\includegraphics[height=3cm]{\pfad/eintraege_01L1.png}

c. 12, False \\
d.

\includegraphics[height=3cm]{\pfad/eintraege_01L2.png}

e. True, ein Fehler wird geworfen
\end{solutionbox}
